\section{Splitting Fields}

\exercise{50.1}{Splitting a quadratic}
The roots of $x^2+3$ are $\pm\sqrt{-3}$. Thus the splitting field is $\Q(\sqrt{-3})$ of degree 2, since the polynomial has no rational root.

\exercise{50.2}{Fourth roots of unity}
The roots of $x^4-1$ are $1,-1,i,-i$. Adjoining $i$ suffices and is necessary. The splitting field $\Q(i)$ has degree 2 over $\Q$.

\exercise{50.3}{Two distinct quadratic radicals}
The splitting field is $\Q(\sqrt2,\sqrt3)$. It has degree 4: if $\sqrt3=a+b\sqrt2$ with rational $a,b$, squaring would give $3=a^2+2b^2+2ab\sqrt2$, hence $ab=0$. The alternatives would require either a rational square root of 3 or a rational square root of $3/2$, both impossible. Thus the second quadratic doubles the first degree.

\exercise{50.4}{Splitting an irreducible cubic}
Let $r=\sqrt[3]3$ and $\omega$ be a primitive cube root of unity. The splitting field is $\Q(r,\omega)$. Eisenstein at 3 gives $[\Q(r):\Q]=3$, and $\omega$ is nonreal while $\Q(r)\subset\R$. Its quadratic polynomial therefore doubles the degree, giving 6.

\exercise{50.5}{Third roots of unity}
Factor $x^3-1=(x-1)(x^2+x+1)$. The quadratic has discriminant $-3$, so its splitting field is $\Q(\sqrt{-3})=\Q(\omega)$, of degree 2.

\exercise{50.6}{Combining square and cube roots}
The roots require $\sqrt2,r=\sqrt[3]2$, and $\omega$. Put $s=2^{1/6}$; then $s=\sqrt2/r$, while $s^2=r$ and $s^3=\sqrt2$. Thus the splitting field is $\Q(s,\omega)$. Eisenstein makes $x^6-2$ irreducible, so $[\Q(s):\Q]=6$. This real field does not contain the nonreal $\omega$, giving total degree 12.

\exercise{50.7}{Automorphisms of the real cubic field}
The order is 1. An automorphism sends the real $r=\sqrt[3]2$ to a root of $x^3-2$ inside the real field $\Q(r)$. Its only real root is $r$, so every automorphism fixes this generator and is the identity.

\exercise{50.8}{Automorphisms of its splitting field}
The splitting field $E=\Q(r,i\sqrt3)$ has degree 6 over $\Q$. Characteristic zero gives six embeddings, all of which are automorphisms because $E$ is a splitting field. Thus the group order is 6.

\exercise{50.9}{Fixing the real cubic generator}
The extension $E/\Q(r)$ has degree 2, obtained by adjoining $i\sqrt3$. Its two automorphisms fix $r$ and send $i\sqrt3$ to itself or its negative. Hence the relative group has order 2.

\exercise{50.10}{A cubic splits in an eight-element field}
The polynomial $f=x^3+x^2+1$ has no root in $\Z_2$, so is irreducible. Thus $\Z_2(\alpha)$ has eight elements. Frobenius preserves roots of polynomials over $\Z_2$, so $\alpha,\alpha^2,\alpha^4$ are roots. They are distinct: a nonzero nonidentity element of the group of seven nonzero elements has order 7, and the exponents 1,2,4 are distinct modulo 7. Hence these exhaust its three roots and $f$ splits there. From $\alpha^3=\alpha^2+1$, the last root is $\alpha^4=\alpha^2+\alpha+1$.

\exercise{50.11}{Correcting the splitting-field criterion}
In the stated condition, ``every polynomial'' must be replaced by ``every irreducible polynomial.'' Otherwise even the splitting field $E=F=\Q$ fails the proposed condition for $(x-1)(x^2-2)$. An algebraic extension $E/F$ is a splitting field of a family of polynomials precisely when each irreducible polynomial over $F$ with one zero in $E$ has all its zeros in $E$.

\exercise{50.12}{Correcting the meaning of splitting}
A polynomial splits over $E$ when it is a nonzero constant times a product of \emph{linear} polynomials in $E[x]$. Merely factoring into lower-degree factors is insufficient: $(x^2-2)(x^2-3)$ factors over $\Q$ but does not split there.

\exercise{50.13}{Bounds for a splitting-field degree}
For a nonconstant degree-$n$ polynomial, $1\le[E:F]\le n!$. Adjoin a first root of degree at most $n$, factor out its linear factor, then adjoin a root of the remaining factor of degree at most $n-1$, and continue. The tower law bounds the product by $n(n-1)\cdots1$. No irreducibility assumption gives a stronger positive lower bound: a polynomial already split over $F$ has degree-one splitting field. For a nonzero constant polynomial the splitting field is $F$ as well.

\exercise{50.14}{Ten statements about splitting fields}
\begin{enumerate}[label=(\alph*)]
\item True. A conjugation isomorphism extends to the closure, and normality preserves $E$; the converse follows by applying an automorphism to a minimal polynomial.
\item False. $\R/\Q$ is not algebraic, whereas a splitting field is algebraic.
\item True: the trivial extension is the splitting field of $x$.
\item True: $\C$ splits $x^2+1$ over $\R$ and is generated by its roots.
\item True: it splits $x^2+1$ over $\Q$.
\item True: $\Q(\pi)$ splits $x^2-\pi^2$ over $\Q(\pi^2)$.
\item False. An isomorphism need not fix $F$ or even have its image in the same field; send the transcendental generator of $E=F=\Q(t)$ to an independent indeterminate.
\item False. For $E=F=\Q(t)$, the embedding $t\mapsto t^2$ has image $\Q(t^2)\subsetneq E$ and lies inside an algebraic closure of $F$.
\item True. This is the defining embedding property of a normal algebraic extension.
\item True. An algebraic closure is generated by all roots of all polynomials over $F$.
\end{enumerate}

\exercise{50.15}{Why irreducibility is essential}
Take $F=E=\Q$, a splitting field of $x$ over $\Q$. The reducible polynomial $(x-1)(x^2-2)$ has the zero 1 in $E$ but its other two zeros are $\pm\sqrt2$, outside $E$. Thus the corollary about an irreducible polynomial having one root in a splitting field does not extend to reducible polynomials.

\exercise{50.16}{Group orders in a tower}
The intended multiplicativity formula has $|G(K/F)|$ on the left; the printed formula repeats $G(E/F)$ there.
\begin{enumerate}[label=(\alph*)]
\item No. For $F=\Q$, $E=\Q(\sqrt[3]2)$ and $K=\Q(\sqrt[3]2,i\sqrt3)$, the group orders are 6,1,2 respectively. Thus $6\ne2\cdot1$.
\item Yes when the relevant fields are splitting fields over the bottom field. Their group orders equal their embedding indices, including for $K/E$, since a normal extension over $F$ is normal over every intermediate field. The multiplicativity of indices then gives $|G(K/F)|=|G(K/E)||G(E/F)|$. This argument does not require separability.
\end{enumerate}

\exercise{50.17}{One polynomial suffices for a finite normal extension}
Choose finitely many generators $\alpha_1,\ldots,\alpha_m$ for $E/F$ and let $p_j$ be their minimal polynomials over $F$. Since $E$ is normal, each $p_j$ splits in $E$. Let $f=\prod_jp_j$. Its splitting field inside $E$ contains every $\alpha_j$, so contains $E$; it is also contained in $E$, since all its roots are there. Hence it equals $E$.

\exercise{50.18}{A quadratic extension is a splitting field}
Choose $\alpha\in E\setminus F$. The tower law gives $E=F(\alpha)$ and a degree-two minimal polynomial $x^2+bx+c$. Its other root, counted with multiplicity, is $-b-\alpha\in E$. Thus the polynomial splits in $E$, and one root generates $E$, proving the claim even in characteristic 2 (where the roots may coincide).

\exercise{50.19}{Containing all conjugates}
If $E/F$ is a splitting field, every irreducible polynomial with a root in $E$ splits there, so all conjugates of every element belong to $E$. Conversely, for each $\alpha\in E$, take its minimal polynomial. The hypothesis places all its roots in $E$. The splitting field of this family is contained in $E$ and contains every $\alpha$, hence is exactly $E$.

\exercise{50.20}{No nontrivial automorphism of the real cubic}
Every automorphism of $\Q(r)$ fixes $\Q$, and the image of $r=\sqrt[3]2$ must satisfy $x^3-2=0$. The other two roots $r\omega,r\omega^2$ are nonreal, whereas the field is real. Thus $r$ is fixed. Since it generates the field, the automorphism is the identity.

\exercise{50.21}{The cyclic subgroup fixing the quadratic radical}
Put $r=\sqrt[3]2$ and $\omega=(-1+i\sqrt3)/2$. The field $\Q(i\sqrt3)=\Q(\omega)$ has degree 2 over $\Q$, so its relative degree in the degree-six splitting field is 3. The map $\rho(r)=r\omega$, $\rho(\omega)=\omega$ is an automorphism, and its powers send $r$ successively to $r\omega^2$ and $r$. It has order 3 and accounts for all three relative embeddings. Thus the group is cyclic, isomorphic to $(\Z_3,+)$.

\exercise{50.22}{The permutation representation on roots}
\begin{enumerate}[label=(\alph*)]
\item For an automorphism \emph{fixing $F$}, $f(\sigma(\alpha))=\sigma(f(\alpha))=0$. Its inverse shows that this gives a permutation of the roots. The base-field qualification is necessary.
\item The roots generate the splitting field over $F$. Agreement on all roots and on $F$ implies agreement on every rational expression in them, hence on the whole field.
\item Restriction to the roots is therefore an injective homomorphism $G(E/F)\hookrightarrow S_\Omega$, where $\Omega$ is the set of distinct roots. Its image is the subgroup of permutations that preserve all algebraic relations among the roots over $F$.
\end{enumerate}

\exercise{50.23}{The full cubic permutation group}
\begin{enumerate}[label=(\alph*)]
\item The field $E=\Q(r,\omega)$ has degree 6 and is separable in characteristic zero. All six rational embeddings are automorphisms, so $|G(E/\Q)|=6$.
\item Its action on the three roots gives an injection into $S_3$. As both groups have six elements, the injection is onto, proving $G(E/\Q)\cong S_3$.
\end{enumerate}

\exercise{50.24}{A prime-root-of-unity extension}
For a prime $p$, every nonidentity $p$th root $\zeta$ is primitive, and its powers give all roots of $x^p-1$. Thus its splitting field is $\Q(\zeta)$. The minimal polynomial $1+x+\cdots+x^{p-1}$ is irreducible over $\Q$ by the prime cyclotomic theorem, so the extension degree is $p-1$ (including degree 1 for $p=2$).

\exercise{50.25}{Uniqueness of splitting fields in different closures}
There is an isomorphism $\tau:\overline F\to\overline F'$ fixing $F$. It maps roots of $f$ to roots of $f$ because the coefficients are fixed. Its inverse gives the reverse inclusion, so it bijects the two full root sets. Since each splitting field is generated by the respective roots, $\tau(E)=E'$. The restriction is the desired $F$-isomorphism.
